% ============================================================
% Mechanistic Validity: A Measurement-Theoretic Framework for
% Evaluating Interpretability Claims
% ============================================================
\pdfoutput=1

% NOTE: This paper is intended for TMLR format.
% If tmlr.sty is available, use:
%   \documentclass[10pt]{article}
%   \usepackage[preprint]{tmlr}
% Otherwise, the following provides compatible formatting.
\documentclass[10pt]{article}

% --- Try TMLR style; fall back to compatible packages ---
\IfFileExists{tmlr.sty}{%
  \usepackage[preprint]{tmlr}%
}{%
  \usepackage[margin=1in]{geometry}%
  \usepackage{times}%
  \usepackage{authblk}%
}

\usepackage[T1]{fontenc}
\usepackage[utf8]{inputenc}
\usepackage{microtype}
\usepackage{amsmath,amssymb}
\usepackage{booktabs}
\usepackage{graphicx}
\usepackage{xcolor}
\usepackage{hyperref}
\usepackage{cleveref}
\usepackage{enumitem}
\usepackage{multirow}
\usepackage{array}
\usepackage{natbib}

% --- Macros ---
\newcommand{\mechval}{\textsc{MechVal}}
\newcommand{\ie}{i.e.}
\newcommand{\eg}{e.g.}
\newcommand{\cf}{cf.}

% tight tables
\newcommand{\tighttable}{\setlength{\tabcolsep}{4pt}\renewcommand{\arraystretch}{1.05}}

\title{Mechanistic Validity: A Measurement-Theoretic Framework for Evaluating Interpretability Claims}

\IfFileExists{tmlr.sty}{}{%
  \author[1]{Elliot Tower}
  \affil[1]{Independent Researcher, \texttt{elliot@elliottower.ai}}
  \date{}
}

\begin{document}

\maketitle

% ============================================================
\begin{abstract}
Mechanistic interpretability (MI) faces a measurement crisis: 99.7\% of sparse autoencoder features are unstable across training runs, attribution patching correlates at $r = 0.006$ with ground truth on synthetic tasks, two of six standard SAE evaluation metrics fail basic diagnostic tests, and 93\% of published MI papers fail execution-grounded reproducibility checks. These findings undermine any claim built on current methodology---including safety-critical claims about deception detection and model monitoring. We present \mechval{}, a measurement-theoretic framework for evaluating MI claims. The framework organizes 27 operational criteria across five validity types (construct, measurement, internal, external, interpretive) in dependency order, assigns claims to one of five verdict tiers (Proposed through Validated), and supports evaluation via 84 metrics spanning five evidence families (causal, structural, behavioral, representational, information-theoretic). The framework is method-agnostic: it applies uniformly to circuits, SAE features, probing classifiers, steering vectors, transcoders, and crosscoders. We describe the framework's theoretical foundations---drawn from philosophy of science, neuroscience, psychometrics, pharmacology, and causal inference---and its open-source implementation, which includes 54 evaluation tasks, 14 calibration gates, and 12 pre-registered mechanistic claim specifications. Application to 13 published circuits demonstrates that the tier system discriminates: claims range from Proposed (Knowledge Neurons, Probing Classifiers, Gender Bias Circuits) through Causally Suggestive (Grokking, Successor Heads, SAE Features) and Mechanistically Supported (IOI Circuit, Greater-Than, Copy Suppression) to Triangulated (Induction Heads), with verdict differences traceable to specific missing evidence.
\end{abstract}

% ============================================================
\section{Introduction}
\label{sec:intro}

When should you trust an interpretability claim? A researcher reports that GPT-2 uses a six-role circuit to solve indirect object identification \citep{wang2023interpretability}, that sparse autoencoder features correspond to human-interpretable concepts \citep{cunningham2023sparse}, or that a linear probe reveals syntactic structure in transformer activations \citep{belinkov2022probing}. Each claim comes with experimental evidence---ablation studies, activation visualizations, accuracy numbers. But how well-supported is the claim, really?

Recent work provides uncomfortable answers. \citet{bricken2024monosemanticity} and \citet{engels2025not} demonstrate that only 197 of 65,536 sparse autoencoder features---0.3\%---are stable across independent training runs. The remaining 99.7\% are training artifacts. Any ``deception feature'' identified using standard SAE methods has a 99.7\% prior probability of being an artifact, not a genuine representation.

The field's primary causal tool fares no better. \citet{kramar2024atp} show that attribution patching---the most widely used method for identifying causally relevant components---correlates at $r = 0.006$ with ground truth on controlled synthetic tasks. This is not a weak correlation; it is no correlation. The method that the field treats as the gold standard for causal attribution provides essentially random rankings of component importance.

The evaluation metrics used to assess these tools are themselves unreliable. The SAEBench audit subjects six standard SAE evaluation metrics to three diagnostic tests: reseed stability, ground-truth correlation, and discriminability. Two of six---Token Prediction Probability (TPP) and Sparse Cross-entropy Recovery (SCR)---fail all three diagnostics. TPP exhibits reseed coefficients of variation between 16\% and 39\%.

Finally, the research itself does not reproduce. MechEvalAgent \citep{bai2026mecheval} performs execution-grounded evaluation of published MI research: 93\% of papers fail reproducibility checks when their code is actually run, and 80\% fail coherence checks between claims and evidence. Narrative review---the kind performed by peer reviewers---missed 51 issues that execution-grounded evaluation caught.

The implication is direct: any claim of the form ``we found a feature that detects deception'' is built on this foundation. If the feature was identified using SAE methods, it is probably a training artifact. If it was validated using attribution patching, the validation is meaningless. If its quality was assessed using TPP or SCR, the assessment is unreliable. If the result was reported in a paper, there is a 93\% chance the code does not reproduce.

This paper presents a framework for addressing this crisis. The \mechval{} framework does not propose new detection methods. It provides a systematic protocol for determining whether any detection method---existing or future---actually works. The framework draws on established validation methodology from five fields (philosophy of science, neuroscience, psychometrics, pharmacology, and causal inference) to define 27 operational criteria, organized across five validity types in dependency order, with five verdict tiers that represent qualitative transitions in evidential status.

The framework is method-agnostic. It evaluates SAE features, neural circuits, probing results, steering vectors, transcoders, and crosscoders using the same criteria. It does not privilege one approach over another. It asks the same questions of each: Is the construct defined? Is the evidence causal? Does the finding generalize? Is the measurement reliable? Is the interpretation justified?

We describe the framework (\Cref{sec:framework}), its benchmark implementation (\Cref{sec:implementation}), its application to 13 published circuits (\Cref{sec:case-studies}), and discuss its implications for the field (\Cref{sec:discussion}).

% ============================================================
\section{The Framework}
\label{sec:framework}

The \mechval{} framework has five interconnected components: validity types (\Cref{sec:validity-types}), criteria (\Cref{sec:criteria}), evidence families (\Cref{sec:evidence-families}), verdict tiers (\Cref{sec:verdict-tiers}), and description modes (\Cref{sec:description-modes}).

\subsection{Five Validity Types in Dependency Order}
\label{sec:validity-types}

The five validity types form a dependency chain: later types cannot be established without first satisfying earlier ones.

\begin{equation}
\text{Construct} \rightarrow \text{Measurement} \rightarrow \text{Internal} \rightarrow \text{External} \rightarrow \text{Interpretive}
\label{eq:dependency}
\end{equation}

Each type addresses a distinct question about the claim:

\paragraph{Construct validity.} Is the thing being measured well-defined? Before any measurement is taken, the construct must be specified precisely enough that the claim is falsifiable, structurally plausible, and distinguishable from neighboring constructs. A ``deception feature'' that cannot be distinguished from an ``uncertainty feature'' fails construct validity regardless of how carefully it is measured. This draws on philosophy of science (Popper's falsifiability, Hempel's operationalism) and psychometrics (construct validity, convergent and discriminant validity).

\paragraph{Measurement validity.} Are the instruments trustworthy? Once the construct is defined, the measurement tools used to detect it must be reliable (stable across repetitions), calibrated (separable from random baselines), and invariant (consistent across conditions). A metric that gives different answers when run with different random seeds is not evidence. This draws on psychometrics (test-retest reliability, measurement invariance) and pharmacology (assay validation).

\paragraph{Internal validity.} Does the evidence support the causal claim? Given a well-defined construct and trustworthy instruments, the evidence must establish that the identified component is causally involved in the behavior---not merely correlated with it. This requires demonstrating necessity (the component is required), sufficiency (the component is enough), and specificity (the component does this task and not everything). This draws on neuroscience (lesion studies, optogenetics, double dissociation) and causal inference \citep{pearl2009causality}.

\paragraph{External validity.} Does the mechanism generalize? Causal evidence on one prompt set, with one ablation method, in one model, establishes internal validity at best. External validity requires that the mechanism generalizes across prompts, ablation methods, related tasks, and ideally across models. This draws on pharmacology (Phase~III generalization, dose-response) and causal inference (transportability; \citealp{pearl2011transportability}).

\paragraph{Interpretive validity.} Is the interpretation of the mechanism correct? Finally, the natural-language description of the mechanism must be justified by the evidence. A head called a ``name mover'' based on direct logit attribution carries theoretical implications beyond what the evidence strictly supports. Interpretive validity requires declaring the description level, matching evidence to that level, and checking for anthropomorphic inflation. This is specific to MI but draws on Marr's levels of analysis \citep{marr1982vision}.

\subsection{Twenty-Seven Criteria}
\label{sec:criteria}

Each validity type is operationalized through specific criteria, totaling 27 across the five types. \Cref{tab:criteria} lists all criteria. Each criterion has a concrete test: a claim either passes, partially passes, or fails each criterion, producing a structured validity profile rather than a single score.

\begin{table}[t]
\centering
\tighttable
\caption{The 27 operational criteria organized by validity type. Each criterion has a concrete test procedure.}
\label{tab:criteria}
\small
\begin{tabular}{@{}llp{6.5cm}@{}}
\toprule
\textbf{Type} & \textbf{ID} & \textbf{Criterion} \\
\midrule
\multirow{5}{*}{Construct} & C1 & Falsifiability --- can the claim be refuted? \\
& C2 & Structural plausibility --- is the mechanism physically possible in the architecture? \\
& C3 & Convergent validity --- do multiple independent methods agree? \\
& C4 & Discriminant validity --- does the measure distinguish this from neighboring constructs? \\
& C5 & Nomological validity --- does the claim fit into a broader theory? \\
\midrule
\multirow{6}{*}{Measurement} & M1 & Reliability --- do repeated measurements give the same answer? \\
& M2 & Baseline separation --- is the score distinguishable from random/untrained baselines? \\
& M3 & Stability --- is the classification robust to perturbation? \\
& M4 & Calibration --- are the numbers meaningful? \\
& M5 & Sensitivity --- can the instrument detect known-true effects? \\
& M6 & Invariance --- does the metric behave consistently across conditions? \\
\midrule
\multirow{5}{*}{Internal} & I1 & Necessity --- is the circuit required for the behavior? \\
& I2 & Sufficiency --- is the circuit enough to produce the behavior? \\
& I3 & Specificity --- does the circuit do this task and not everything? \\
& I4 & Double dissociation --- can this circuit be separated from other circuits? \\
& I5 & Confound control --- are alternative explanations ruled out? \\
\midrule
\multirow{6}{*}{External} & E1 & Intervention reach --- do different intervention methods agree? \\
& E2 & Prompt generalization --- does it work on diverse prompts? \\
& E3 & Cross-task generalization --- does the mechanism transfer to related tasks? \\
& E4 & Cross-model generalization --- does the mechanism appear in other models? \\
& E5 & Graded response --- does partial ablation produce partial effects? \\
& E6 & Novel prediction --- does the mechanism predict new, untested behaviors? \\
\midrule
\multirow{5}{*}{Interpretive} & V1 & Level declaration --- at what description mode is the claim made? \\
& V2 & Level-evidence match --- does the evidence support claims at that level? \\
& V3 & Alternative level --- could the evidence be explained at a different level? \\
& V4 & Anthropomorphism check --- is the interpretation projecting human concepts? \\
& V5 & Scope declaration --- what does the claim explicitly not cover? \\
\bottomrule
\end{tabular}
\end{table}

\subsection{Evidence Families}
\label{sec:evidence-families}

Metrics are organized into six evidence families, each providing a different kind of evidence about a mechanistic claim. Each metric belongs to exactly one family.

\begin{itemize}[nosep]
    \item \textbf{Causal (A):} Effects of interventions---ablation, patching, clamping, resampling. Examples: activation patching, DAS-IIA, CATE, mediation analysis, sigma ablation.
    \item \textbf{Structural (B):} Weight-space properties, requiring no forward pass. Examples: SVD, effective rank, OV/QK composition norms, template distance.
    \item \textbf{Information-theoretic (C):} Information flow and dependencies. Examples: mutual information, partial information decomposition, OCSE, transfer entropy.
    \item \textbf{Behavioral (D):} Input-output behavior under manipulation. Examples: faithfulness, logit difference, KL divergence, dose-response curves.
    \item \textbf{Representational (E):} Geometry and content of internal representations. Examples: CKA, RSA, linear probes, attention entropy.
    \item \textbf{Measurement (F):} Meta-properties of other metrics---calibrations that assess whether the primary metrics are trustworthy. Examples: bootstrap stability, seed variance, convergent validity, discriminant validity.
\end{itemize}

The distinction between families matters because convergent evidence---the same conclusion supported by metrics from different families---is qualitatively stronger than redundant evidence from the same family. An ablation study (causal) that agrees with a weight-space analysis (structural) provides more support than two ablation studies using different ablation methods. The framework tracks evidence family coverage explicitly.

\subsection{Verdict Tiers}
\label{sec:verdict-tiers}

Claims are assigned to one of five verdict tiers representing qualitative transitions in evidential status. Two additional labels handle special cases. The tiers form a strict hierarchy: each requires all evidence from the tier below plus additional evidence.

\begin{table}[t]
\centering
\tighttable
\caption{Verdict tiers with requirements. Each tier subsumes the requirements of all lower tiers.}
\label{tab:verdicts}
\small
\begin{tabular}{@{}lp{4.5cm}p{5.5cm}@{}}
\toprule
\textbf{Tier} & \textbf{Meaning} & \textbf{Requirements beyond previous tier} \\
\midrule
Proposed & Structural or representational evidence only & Construct validity criteria met (C1--C5) \\
Causally Suggestive & Necessity shown, sufficiency not established & Internal validity I1 (necessity via ablation) \\
Mechanistically Supported & Necessity + sufficiency with consistent methods & I1 + I2 (sufficiency), intervention reach (E1), at least 2 ablation variants \\
Triangulated & Multiple converging lines of independent evidence & Full internal + external + construct criteria from $\geq 3$ evidence families \\
Validated & All five validity types addressed & Measurement + interpretive validity with explicit baselines \\
\midrule
Underdetermined & Evidence consistent with multiple mechanisms & Cannot resolve between rival specifications \\
Disconfirmed & Fails decisively on a key criterion & Negative result \\
\bottomrule
\end{tabular}
\end{table}

The tier system is designed to be informative at every level. A claim at \emph{Proposed} is not invalid---it is a claim whose evidential status is precisely characterized. The verdict tells researchers exactly what additional evidence would advance the claim: a Proposed claim needs causal intervention (I1) to reach Causally Suggestive; a Causally Suggestive claim needs sufficiency evidence (I2) and method consistency (E1) to reach Mechanistically Supported.

\subsection{Description Modes}
\label{sec:description-modes}

MI claims operate at different levels of description, following Marr's hierarchy but adapted for transformer circuits. The framework defines seven description modes forming a partial order by evidential commitment:

\begin{enumerate}[nosep]
    \item \textbf{Computational:} What function does the system compute, and why?
    \item \textbf{Algorithmic:} What procedure does it execute?
    \item \textbf{Representational:} What information is encoded, and in what geometry?
    \item \textbf{Implementational-Functional:} What input-output transformation does each component perform?
    \item \textbf{Implementational-Connectomic:} How are components wired?
    \item \textbf{Implementational-Topographic:} Which components are involved?
    \item \textbf{Implementational-Statistical:} What do activations look like? (Orthogonal to the others.)
\end{enumerate}

Higher modes require all evidence from lower modes plus bridging evidence that connects levels. A claim at the algorithmic level (``the model runs a detect-inhibit-copy algorithm'') requires both the topographic evidence (which heads are involved) and the additional evidence that the heads implement the claimed algorithm. Mode tags are applied to verdicts, not to individual metrics---a single experiment can provide evidence for multiple modes simultaneously.

% ============================================================
\section{Theoretical Foundations}
\label{sec:foundations}

The framework synthesizes validation methodology from five established fields. This section describes the specific contributions of each.

\subsection{Philosophy of Science}
\label{sec:phil-sci}

The construct validity criteria (C1--C5) draw on Popper's falsifiability criterion, Mayo's severe testing \citep{mayo2018statistical}, and Glennan's new mechanical philosophy \citep{glennan2017new}. C1 requires that a mechanistic claim state what evidence would disprove it---a requirement routinely violated in MI, where circuit labels are often applied post-hoc without pre-registered predictions. C5 (nomological validity) requires that the claim fit into a broader theoretical network, following the nomological network concept from Cronbach and Meehl.

The verdict tier system is modeled on the hierarchy of evidential support in philosophy of science: observation (Proposed), single-method causal evidence (Causally Suggestive), multi-method causal evidence (Mechanistically Supported), convergent evidence from independent methods (Triangulated), and comprehensive evaluation (Validated).

\subsection{Neuroscience}
\label{sec:neuro}

The internal validity criteria (I1--I5) are adapted from neuroscience's toolkit for establishing that a neural circuit implements a computation. I1 (necessity) corresponds to lesion studies: removing the circuit disrupts the behavior. I2 (sufficiency) corresponds to stimulation: activating the circuit alone produces the behavior. I4 (double dissociation) is the gold standard in neuropsychology for establishing functional specificity---showing that circuit A is necessary for task X but not task Y, while circuit B is necessary for task Y but not task X.

The critical distinction between necessity and sufficiency---well-established in neuroscience but often elided in MI---is central to the framework. Most published circuits demonstrate necessity (ablation degrades behavior) but not sufficiency (the circuit alone can produce the behavior). The framework treats this distinction as a qualitative transition between verdict tiers.

\subsection{Psychometrics}
\label{sec:psychometrics}

The measurement validity criteria (M1--M6) draw on psychometric principles of test construction. M1 (reliability) is test-retest reliability: the same metric applied to the same model should give the same answer. M6 (invariance) is measurement invariance: the metric should behave consistently across conditions (\eg, different prompt sets, different ablation methods). The calibration gates (F01--F15 in the implementation) operationalize these principles as concrete statistical tests.

The distinction between convergent and discriminant validity (C3 and C4) comes directly from Campbell and Fiske's multitrait-multimethod matrix. A valid construct should show convergent validity (multiple methods agree) and discriminant validity (the construct is distinguishable from related constructs). In MI terms: a ``deception feature'' should be detected by both probing and activation patching (convergent), and should not also fire on uncertainty, hedging, or politeness (discriminant).

\subsection{Pharmacology}
\label{sec:pharma}

The external validity criteria, particularly E5 (graded response), draw on pharmacology's dose-response methodology. A causally relevant component should produce graded effects: partial ablation should produce partial behavioral change, and the relationship should be monotonic or at least systematic. The framework treats dose-response as a minimum bar for causal claims about magnitude---a requirement absent from most MI evaluations, which typically report binary ablation (all-or-nothing).

The three-phase protocol used in the case studies (measurement reliability, causal validity, generalization) mirrors the phase structure of clinical trials: Phase~I establishes safety and dosing (analogous to measurement calibration), Phase~II establishes efficacy (causal validity), and Phase~III establishes generalization.

\subsection{Causal Inference}
\label{sec:causal}

The benchmark architecture (\Cref{sec:implementation}) is grounded in causal inference methodology. Track~1 (Circuit Localization) corresponds to causal discovery \citep{spirtes2000causation}---identifying the causal structure from data. Track~2 (Causal Variable Localization) corresponds to causal abstraction \citep{geiger2023causal}---finding where high-level causal variables are represented. Track~3 (Causal Model Testing) corresponds to model testing in the Spirtes-Glymour-Scheines and Pearl frameworks---testing whether a proposed causal model is consistent with interventional data.

The four scoring views map to established causal inference concepts: V1 (Effect Estimation) to Pearl/Rubin effect estimation, V2 (Transportability) to Pearl/Bareinboim transportability theory \citep{pearl2011transportability}, V3 (Counterfactual Verification) to Pearl's rung-3 counterfactuals, and V4 (Mechanism Adjudication) to SEM equivalent-models testing.

\Cref{tab:crossfield} shows the complete cross-field mapping.

\begin{table}[t]
\centering
\tighttable
\caption{Cross-field mapping of framework components to established methodology.}
\label{tab:crossfield}
\small
\begin{tabular}{@{}lp{2.2cm}p{2.2cm}p{2.0cm}p{1.8cm}p{2.0cm}@{}}
\toprule
\textbf{Component} & \textbf{Causal Inf.} & \textbf{Neuroscience} & \textbf{Psychometrics} & \textbf{Pharmacology} & \textbf{Phil.\ of Sci.} \\
\midrule
Track 1 & Causal discovery & Circuit mapping & Exploratory FA & Target discovery & Abduction \\
Track 2 & Causal abstraction & Functional localization & Construct oper. & Target charac. & Natural kinds \\
Track 3 & Model testing & Circuit dissection & Confirmatory FA & Proof of mech. & H-D testing \\
V1 & ATE/CATE & Perturbation quant. & Effect size est. & Dose-response & Parameter est. \\
V2 & Transportability & Cross-condition gen. & Meas.\ invariance & Phase III gen. & Novel prediction \\
V3 & Rung-3 counterfact. & Optogenetics & --- & --- & Counterfactual \\
V4 & Equiv.\ models & --- & Model comparison & --- & Crucial expt. \\
\bottomrule
\end{tabular}
\end{table}

% ============================================================
\section{Implementation}
\label{sec:implementation}

The framework is implemented as an open-source Python package. This section describes the benchmark architecture, the core data structures, and the current scale of the implementation.

\subsection{Three Tracks, Four Views, Four Gates}
\label{sec:3t4v4g}

The benchmark is organized into three competitive tracks, four scoring views, and four precondition gates.

\paragraph{Track 1: Circuit Localization (Causal Discovery).} Given a model and task, discover which components are causally relevant. Input: model + task. Output: edge set / circuit. Score: Circuit Model Distance (CMD), faithfulness. Methods compete on the same input to find the best circuit.

\paragraph{Track 2: Causal Variable Localization (Causal Abstraction).} Given a model, task, and hypothesized high-level variable, find where that variable is represented. Input: model + task + variable. Output: aligned subspace. Score: Interchange Intervention Accuracy (IIA).

\paragraph{Track 3: Causal Model Testing.} Given a model, task, and pre-registered mechanistic hypothesis (a \texttt{MechanisticClaimSpec}), test whether the model's interventional behavior matches the hypothesis's predictions. Input: model + task + claim specification. Output: verdict profile + claim ceiling. Score: confirmation rate, negative control rate. This track is the framework's primary original contribution---theory-driven hypothesis testing rather than blind discovery.

The key distinction between Tracks 1 and 3 mirrors the distinction between exploratory and confirmatory factor analysis in psychometrics. Track~1 is unsupervised: discover whatever circuit exists. Track~3 is supervised: test whether a specific proposed mechanism is actually implemented. Methods compete on how well their candidate circuits satisfy pre-registered predictions.

\paragraph{Views} are scoring aggregations over existing metrics---different causal-inference-grounded lenses on the same results. They do not introduce new measurement infrastructure:
\begin{itemize}[nosep]
    \item V1 (Causal Effect Estimation): How large is each component's contribution?
    \item V2 (Causal Transportability): Does the mechanism generalize?
    \item V3 (Counterfactual Verification): Do counterfactual interventions produce expected results?
    \item V4 (Mechanism Adjudication): When multiple mechanisms could explain the data, which is better supported?
\end{itemize}

\paragraph{Gates} are binary preconditions that must pass before tracks and views produce meaningful results:
\begin{itemize}[nosep]
    \item G0 (Construct Operationalization): Is the construct defined independently of the instruments?
    \item G1 (Measurement Calibration): Are the metrics stable across seeds and separable from baselines?
    \item G2 (Causal Identifiability): Can the specified causal effects be estimated with available interventions?
    \item G3 (Confound/Superposition Risk): Are interventions confounded by polysemantic collateral damage?
\end{itemize}

G1 is particularly important. An IIA score of 0.48 without knowing the random-vector baseline is not evidence. The gate requires bootstrap stability (scores stable under resampling), seed variance (stable across random seeds), and baseline separation (distinguishable from random and untrained-model baselines) before any downstream result is considered meaningful.

\subsection{MechanisticClaimSpec}
\label{sec:claimspec}

Track~3 requires a pre-registered mechanistic hypothesis formalized as a \texttt{MechanisticClaimSpec}---a Pydantic v2 data model that encodes:

\begin{itemize}[nosep]
    \item \textbf{Computational steps:} Nodes in the mechanism DAG, each with a name, category, description, and mapping to specific model components (attention heads, MLPs, neurons, SAE features).
    \item \textbf{Edges:} Directed connections between steps, specifying the communication mechanism (residual stream, attention composition, etc.).
    \item \textbf{Positive predictions:} Testable claims about what \emph{should} happen under intervention (\eg, ``ablating duplicate-token-detection heads reduces logit difference by $\geq 0.2$'').
    \item \textbf{Negative controls:} Claims about what should \emph{not} happen (\eg, ``ablating name-mover heads does not affect upstream S-inhibition heads'').
    \item \textbf{Rival specifications:} IDs of competing mechanistic hypotheses.
    \item \textbf{Gate metadata:} Identifiability (G2) and superposition risk (G3) assessments.
\end{itemize}

When \texttt{verify(spec)} runs, it executes each prediction by routing to the appropriate metric (typically role ablation---ablating all heads in one functional role and measuring the effect), compares measured values against expected directions and thresholds, and aggregates results into a confirmation rate, negative control rate, claim ceiling (highest description mode with all predictions passing), and verdict tier.

The pre-registration requirement is critical. It prevents post-hoc rationalization---the practice of discovering a circuit and then constructing a narrative that fits the observations without ever testing predictions the narrative makes. By requiring predictions and negative controls \emph{before} verification, the framework enforces the confirmatory discipline that the field currently lacks.

\subsection{Scale of Implementation}
\label{sec:scale}

\Cref{tab:scale} summarizes the current implementation.

\begin{table}[t]
\centering
\tighttable
\caption{Current scale of the \mechval{} implementation.}
\label{tab:scale}
\small
\begin{tabular}{@{}lr@{}}
\toprule
\textbf{Component} & \textbf{Count} \\
\midrule
Tasks & 54 (12 full circuit, 10 proxy, 25 generator-only, 7 planned) \\
Metrics & 84 across 5 evidence families \\
Calibrations & 14 \\
Claim specifications & 12 pre-registered \\
Linguistic domains & 10 \\
Worked case studies & 13 published circuits evaluated \\
Architecture decisions & 14 ADRs \\
\bottomrule
\end{tabular}
\end{table}

The 84 metrics span all five evidence families: 32 causal (including activation patching, DAS-IIA, mediation, CATE, sigma ablation, causal scrubbing, EAP, path patching), 18 structural (composition scores, template distance, graph topology, spectral analysis), 9 information-theoretic (PID, OCSE, NOTEARS, mutual information, transfer entropy), 20 behavioral (dose-response, faithfulness, generalization gap, boundary analysis), and 5 representational (CKA, probes, attention entropy).

The 14 calibrations operationalize measurement validity: bootstrap stability (F01), seed variance (F02), convergent validity (F03), discriminant validity (F04), internal consistency (F05), split-half and test-retest reliability (F06), measurement invariance (F07), distributional characterization (F09), distributional stability (F10), nomological validity (F11), incremental validity (F12), ablation invariance (F13), method invariance (F14), and certified stability (F15). Three are hard gates: F01 (bootstrap), F09 (random baseline), and F10 (untrained baseline) must pass before downstream results are meaningful.

The 12 claim specifications cover circuits of varying complexity: IOI (6 steps, 15 heads, 7 predictions, 5 negative controls), Greater-Than (3 steps, 7 heads), Induction (2 steps, 7 heads), SVA (4 steps, 12 heads), Copy Suppression (3 steps, 7 heads), Gendered Pronoun (3 steps, 5 heads), Acronym (3 steps, 8 heads), RTI (4 steps, 15 heads), and four epistemic framing variants discovered via different methods (manual, activation patching, EAP, and broad scan).

The 54 tasks span 10 linguistic domains: agreement, binding, coreference, morphology, phonology, pragmatics, semantics, syntax, math, and patterns. Tasks are registered in a Gymnasium-style registry and loaded by ID.

\subsection{API}
\label{sec:api}

The framework provides both a Python API and a CLI:

\begin{verbatim}
import mechanistic_validity as mv

# Load task and run metric
task = mv.load_task("ioi")
results = mv.run("activation_patching", tasks=["ioi"])

# Track 3: verify a claim specification
spec = task.get_claim_spec()
result = mv.verify(spec, device="cpu")
print(result.confirmation_rate)  # 0.0 to 1.0
print(result.verdict_tier)       # VerdictTier enum

# Calibrate a metric
mv.calibrate("bootstrap", tasks=["ioi"])

# Check precondition gates
mv.check_gate("measurement_calibration", task="ioi")
\end{verbatim}

The \texttt{ArtifactAdapter} pattern allows the framework to evaluate learned representation artifacts (SAEs, factor banks, transcoders, crosscoders) without coupling to any specific implementation. Adapters expose standard methods (\texttt{directions}, \texttt{activations}, \texttt{ablate}) regardless of the underlying architecture.

% ============================================================
\section{Case Studies: Thirteen Published Circuits}
\label{sec:case-studies}

We apply the framework to 13 published MI results, evaluating each through all five validity lenses. The case studies are not intended to rank papers---they demonstrate what the framework looks like in practice and show that the tier system discriminates between well-validated and poorly-validated claims.

\subsection{Verdict Assignments}

\Cref{tab:case-studies} presents the verdict assignments. The key finding is that the tier system produces meaningful discrimination: the 13 claims span the full range from Proposed to Triangulated, with verdict differences traceable to specific missing evidence rather than arbitrary scoring.

\begin{table}[t]
\centering
\tighttable
\caption{Verdict assignments for 13 published circuits. Verdicts are based on published evidence as of evaluation date.}
\label{tab:case-studies}
\small
\begin{tabular}{@{}llp{6.5cm}@{}}
\toprule
\textbf{Verdict} & \textbf{Circuit} & \textbf{Key limiting criterion} \\
\midrule
\multirow{1}{*}{Triangulated} & Induction Heads & Simple mechanism, broad replication, thick nomological network \\
\midrule
\multirow{3}{*}{\parbox{2.3cm}{Mechanistically\\Supported}} & IOI Circuit & Method-conditional faithfulness (87\% mean ablation, $<$50\% other methods); specificity untested \\
& Greater-Than & Strong structural plausibility (C2); limited prompt generalization \\
& Copy Suppression & Unusually clean specificity; sufficiency partially established \\
\midrule
\multirow{5}{*}{\parbox{2.3cm}{Causally\\Suggestive}} & Grokking & Validated within toy scope; toy-to-real gap is the limitation \\
& Successor Heads & Cross-domain generalization as convergent evidence; ablation method-conditional \\
& Docstring Circuit & Label ambiguity: ``variable binding'' vs.\ simpler ``positional copying'' \\
& SAE Features & Thin nomological network; features may be dictionary properties, not model properties \\
& Othello World Model & ``World model'' label exceeds evidence (``linearly decodable''); interpretive inflation \\
\midrule
\multirow{4}{*}{Proposed} & Knowledge Neurons & Strong intervention but weak mechanistic story \\
& Superposition & Validated theory in toy models; awaits real-model confirmation \\
& Probing Classifiers & Decodability without intervention = no internal validity \\
& Gender Bias Circuits & Construct incoherence: bias and knowledge share circuits \\
\bottomrule
\end{tabular}
\end{table}

\subsection{Illustrative Examples}

We describe three case studies in detail to illustrate how the framework produces its verdicts.

\paragraph{Induction Heads (Triangulated).} \citet{olsson2022context} identify a two-head composition mechanism for in-context copying: previous-token heads attend to the token before a repeated sequence, and induction heads use this signal to predict the next token. The claim reaches Triangulated because it satisfies criteria across multiple independent evidence families. The mechanism has been replicated across model scales and architectures (E4), confirmed through both activation-level and weight-level analysis (C3 convergent validity from different families), and produces clean dose-response effects (E5). The nomological network is thick: the mechanism connects to in-context learning theory, training dynamics, and cross-model structural predictions. The relative simplicity of the mechanism (two functional roles, clear compositional structure) makes each criterion easier to satisfy.

\paragraph{IOI Circuit (Mechanistically Supported).} \citet{wang2023interpretability} identify 26 attention heads forming a six-role circuit for indirect object identification. The circuit demonstrates strong necessity (I1) and sufficiency (I2): running the model with everything outside the circuit mean-ablated recovers 87\% of the full model's logit difference. However, \citet{miller2024transformer} show that this 87\% figure is specific to mean ablation; under other ablation methods, faithfulness drops below 50\%. This method-conditional result partially satisfies E1 (intervention reach) but reveals that the causal claim is weaker than the headline number suggests. Specificity (I3) is untested: the authors do not report whether ablating the IOI circuit degrades unrelated tasks. A formal double-dissociation has not been conducted. The circuit reaches Mechanistically Supported rather than Triangulated because the evidence comes primarily from one methodological family (ablation-based) and key criteria (I3, I4, I5, E4) remain unaddressed.

\paragraph{Probing Classifiers (Proposed).} Linear probing \citep{belinkov2022probing} trains a classifier on activations to test whether a concept is linearly decodable. The core limitation is fundamental: \emph{a measurement without an intervention is a measurement without internal validity}. Probe success establishes that information is decodable from the activation space, but not that the model uses that information during inference. High probe accuracy is consistent with genuine encoding, incidental linear separability in high-dimensional space, and confound encoding (a correlated feature rather than the target). Without causal follow-up (DAS, intervention along the probe direction), the claim cannot advance beyond Proposed. Probing \emph{with} causal follow-up can reach Causally Suggestive; with additional control tasks and cross-method convergence, it can reach Mechanistically Supported. The framework does not dismiss probing---it precisely characterizes what probing alone does and does not establish.

\subsection{Patterns Across Case Studies}

Several patterns emerge from evaluating these claims side by side.

\paragraph{The sufficiency gap.} Most circuits demonstrate necessity (I1) but not sufficiency (I2). The I1$\rightarrow$I2 transition---from ``removing this breaks the behavior'' to ``this alone produces the behavior''---is the most common barrier between Causally Suggestive and Mechanistically Supported.

\paragraph{Method-conditional results.} The IOI circuit's headline numbers are specific to mean ablation. Ablation type is part of the causal claim, not an implementation detail. The framework captures this through E1 (intervention reach), which requires agreement across ablation methods.

\paragraph{The toy-model ceiling.} Grokking and Superposition achieve high validity within toy scope---but the gap between toy-model proof-of-concept and real-model confirmation remains the field's central challenge. The framework handles this through the scope declaration criterion (V5): a claim Validated within a declared scope is a legitimate result, not a failure.

\paragraph{Interpretive inflation.} Labels like ``world model,'' ``knowledge neuron,'' and ``deception feature'' carry theoretical implications beyond what the evidence supports. The framework systematically identifies where labels exceed evidence through V4 (anthropomorphism check) and V5 (scope declaration).

\paragraph{Construct incoherence.} Gender Bias Circuits fail not because evidence is lacking but because the construct itself---``bias'' separable from ``legitimate gender processing''---may not be well-posed. Sometimes the right answer is ``this question is not well-defined,'' not ``we need more data.'' The framework captures this through the Construct validity criteria, which must be satisfied before any measurement is meaningful.

% ============================================================
\section{Discussion}
\label{sec:discussion}

\subsection{Method Agnosticism}

The framework evaluates any MI claim type using the same criteria. It applies to:

\begin{itemize}[nosep]
    \item \textbf{Circuits:} Sets of attention heads, MLPs, and edges implementing a computation.
    \item \textbf{SAE features:} Directions learned by sparse autoencoders.
    \item \textbf{Probing classifiers:} Linear classifiers trained on activations.
    \item \textbf{Steering vectors:} Directions used for activation addition/subtraction.
    \item \textbf{Transcoders:} Cross-layer feature extractors.
    \item \textbf{Crosscoders:} Multi-model feature extractors.
\end{itemize}

This agnosticism is not a design choice---it is a consequence of grounding the framework in validity theory rather than in the specifics of any detection method. The questions ``Is the construct well-defined?'' and ``Is the evidence causal?'' apply regardless of whether the evidence comes from ablation patching, dictionary learning, or linear probing.

\subsection{Relationship to Existing Benchmarks}

The Mechanistic Interpretability Benchmark (MIB; \citealp{mueller2025mib}) provides competitive benchmarks for circuit discovery. \mechval{} is complementary: MIB evaluates how well a method discovers circuits (Track~1); \mechval{} evaluates how well-validated the discovered circuit is (Track~3). A circuit that achieves high CMD on MIB may still be poorly validated if it lacks sufficiency evidence, dose-response testing, or cross-method consistency. The two frameworks can be composed: discover a circuit via MIB-style methods, then evaluate it via \mechval{}'s validity criteria.

SAEBench evaluates SAE quality via standardized metrics. \mechval{}'s measurement validity criteria (M1--M6) and calibration gates (F01--F15) address a prior question: are the SAEBench metrics themselves trustworthy? The SAEBench audit's finding that TPP and SCR fail basic diagnostics is precisely the kind of measurement validity failure that the framework's gates are designed to catch.

\subsection{Limitations}

\paragraph{Scope.} The current implementation is applied to GPT-2 Small. Extension to larger models requires substantially more compute and may reveal that some criteria (\eg, E4 cross-model generalization) are harder to satisfy than the small-model case studies suggest.

\paragraph{Verdict assignment subjectivity.} While the 27 criteria are operationally defined, the aggregation from criterion-level verdicts (pass/partial/fail) to tier-level verdicts involves judgment. Two evaluators could assign different tiers to the same circuit if they weight criteria differently. The framework mitigates this by requiring criterion-level transparency: the structured validity profile is the primary output, and the tier is a summary.

\paragraph{No novel experiments.} This work is a theoretical contribution. The case study verdicts are based on published evidence, not new experiments. Running the framework's verification pipeline on all 12 claim specifications with actual model interventions is future work that requires GPU compute.

\paragraph{Coverage gaps.} Of the 54 tasks, only 12 have full circuit definitions with claim specifications. The remaining 42 have prompt generators but lack verified circuits. The framework's utility increases with circuit coverage, which is bottlenecked by the labor-intensive process of circuit discovery and specification.

\subsection{Implications for AI Safety}

The measurement crisis described in \Cref{sec:intro} has direct implications for AI safety. Organizations are beginning to deploy internal monitors based on MI findings---probes that flag potential deception, steering vectors that suppress harmful outputs, circuit-level classifiers that trigger human review. If those monitors are built on unvalidated features, they provide false confidence: the system appears monitored when it is not.

The framework provides a concrete protocol for assessing monitor trustworthiness. A ``deception feature'' can be evaluated through the same criteria as any other MI claim: Is it well-defined (C1--C5)? Is it reliably measurable (M1--M6)? Is it causally necessary and sufficient (I1--I2)? Does it generalize across prompts and models (E1--E6)? Does the ``deception'' label match the evidence (V1--V5)?

The verdict tier system provides deployment guidance. A feature at Proposed should not be deployed as a monitor. A feature at Causally Suggestive provides preliminary evidence but lacks sufficiency testing. A feature at Mechanistically Supported or above has the minimum evidential basis for cautious deployment with ongoing validation. This is not a guarantee---it is a calibrated assessment of evidential status, which is what deployment decisions require.

% ============================================================
\section{Conclusion}
\label{sec:conclusion}

The mechanistic interpretability field has made remarkable progress in identifying circuits, features, and representations in neural networks. But progress in discovery has outpaced progress in validation. The result is a field that generates claims faster than it can evaluate them, using tools whose reliability is uncertain, producing results that do not reproduce.

The \mechval{} framework addresses this gap by providing the measurement infrastructure that the field currently lacks. Its contribution is not a new detection method but a systematic protocol for determining whether any detection method actually works. The 27 criteria, five validity types, five evidence families, and five verdict tiers provide a common language for evaluating MI claims---whether they concern circuits, SAE features, probes, or steering vectors.

The framework's application to 13 published circuits demonstrates that the tier system discriminates meaningfully. The discrimination is informative: verdict differences are traceable to specific missing evidence, and each verdict tells researchers exactly what additional evidence would advance the claim. The framework does not rank papers or methods. It characterizes the evidential status of claims, making the gaps visible and the path forward explicit.

The framework is open-source and method-agnostic. Its contribution is a public good: the measurement infrastructure that makes interpretability claims trustworthy.

% ============================================================
\section*{Acknowledgments}

The theoretical framework draws on insights from the mechanistic interpretability community, particularly the open problems identified by \citet{sharkey2026open} and the actionability framework proposed by \citet{orgad2026actionable}. The cross-field methodology benefits from established validation frameworks in philosophy of science \citep{mayo2018statistical}, neuroscience, psychometrics, pharmacology, and causal inference \citep{pearl2009causality, spirtes2000causation}.

% ============================================================
\bibliographystyle{plainnat}

\begin{thebibliography}{99}

\bibitem[Bai et~al.(2026)]{bai2026mecheval}
Bai, X., Baumgartner, T., Sun, C., Holtzman, A., and Tan, C.
\newblock MechEvalAgent: Execution-grounded evaluation of mechanistic interpretability research.
\newblock \emph{Preprint}, 2026.

\bibitem[Belinkov(2022)]{belinkov2022probing}
Belinkov, Y.
\newblock Probing classifiers: Promises, shortcomings, and advances.
\newblock \emph{Computational Linguistics}, 48(1):207--219, 2022.

\bibitem[Bricken et~al.(2024)]{bricken2024monosemanticity}
Bricken, T., Templeton, A., Batson, J., Chen, B., Jermyn, A., Conerly, T., Turner, N., Anil, C., Denison, C., Askell, A., et~al.
\newblock Towards monosemanticity: Decomposing language models with dictionary learning.
\newblock \emph{Transformer Circuits Thread}, 2024.

\bibitem[Cunningham et~al.(2023)]{cunningham2023sparse}
Cunningham, H., Ewart, A., Riggs, L., Huben, R., and Sharkey, L.
\newblock Sparse autoencoders find highly interpretable features in language models.
\newblock \emph{arXiv preprint arXiv:2309.08600}, 2023.

\bibitem[Engels et~al.(2025)]{engels2025not}
Engels, J., Liao, I., Michaud, E.~J., Gurnee, W., and Tegmark, M.
\newblock Not all language model features are linear.
\newblock In \emph{NeurIPS}, 2025.

\bibitem[Garcia-Carrasco et~al.(2024)]{garcia2024acronym}
Garcia-Carrasco, J., et~al.
\newblock Acronym prediction with ACDC.
\newblock In \emph{AISTATS}, 2024.

\bibitem[Geiger et~al.(2023)]{geiger2023causal}
Geiger, A., Wu, Z., Lu, H., Rozner, J., Kreiss, E., Icard, T., Goodman, N.~D., and Potts, C.
\newblock Causal abstraction: A theoretical foundation for mechanistic interpretability.
\newblock \emph{JMLR}, 2023.

\bibitem[Glennan(2017)]{glennan2017new}
Glennan, S.
\newblock \emph{The New Mechanical Philosophy}.
\newblock Oxford University Press, 2017.

\bibitem[Hanna et~al.(2023)]{hanna2023greater}
Hanna, M., Liu, O., and Variengien, A.
\newblock How does {GPT}-2 compute greater-than over the calendar?
\newblock In \emph{NeurIPS}, 2023.

\bibitem[Kr{\'a}mar et~al.(2024)]{kramar2024atp}
Kr{\'a}mar, J., Lieberum, T., Shah, R., and Neel, N.
\newblock Relevance propagation for transformers.
\newblock In \emph{NeurIPS}, 2025.

\bibitem[Lazo et~al.(2025)]{lazo2025sva}
Lazo, P., et~al.
\newblock Subject-verb agreement circuits in {GPT}-2.
\newblock \emph{arXiv preprint arXiv:2506.22105}, 2025.

\bibitem[Marr(1982)]{marr1982vision}
Marr, D.
\newblock \emph{Vision: A Computational Investigation into the Human Representation and Processing of Visual Information}.
\newblock MIT Press, 1982.

\bibitem[Mayo(2018)]{mayo2018statistical}
Mayo, D.~G.
\newblock \emph{Statistical Inference as Severe Testing: How to Get Beyond the Statistics Wars}.
\newblock Cambridge University Press, 2018.

\bibitem[McDougall et~al.(2023)]{mcdougall2023copy}
McDougall, C., Conmy, A., Rushing, C., McGrath, T., and Neel, N.
\newblock Copy suppression: Comprehensively understanding an attention head.
\newblock \emph{arXiv preprint arXiv:2310.04625}, 2023.

\bibitem[Miller et~al.(2024)]{miller2024transformer}
Miller, J., et~al.
\newblock Transformer circuit faithfulness metrics are not robust.
\newblock \emph{arXiv preprint arXiv:2407.08734}, 2024.

\bibitem[Mueller et~al.(2025)]{mueller2025mib}
Mueller, A., et~al.
\newblock {MIB}: Mechanistic Interpretability Benchmark.
\newblock \emph{arXiv preprint arXiv:2504.13151}, 2025.

\bibitem[Olsson et~al.(2022)]{olsson2022context}
Olsson, C., Elhage, N., Nanda, N., Joseph, N., DasSarma, N., Henighan, T., Mann, B., Askell, A., Bai, Y., Chen, A., et~al.
\newblock In-context learning and induction heads.
\newblock \emph{arXiv preprint arXiv:2209.11895}, 2022.

\bibitem[Orgad et~al.(2026)]{orgad2026actionable}
Orgad, H., Barez, F., et~al.
\newblock Interpretability can be actionable.
\newblock In \emph{ICML}, 2026.

\bibitem[Pearl(2009)]{pearl2009causality}
Pearl, J.
\newblock \emph{Causality: Models, Reasoning, and Inference}.
\newblock Cambridge University Press, 2nd edition, 2009.

\bibitem[Pearl and Bareinboim(2011)]{pearl2011transportability}
Pearl, J. and Bareinboim, E.
\newblock Transportability of causal and statistical relations: A formal approach.
\newblock In \emph{AAAI}, 2011.

\bibitem[Sharkey et~al.(2026)]{sharkey2026open}
Sharkey, L., et~al.
\newblock Open problems in mechanistic interpretability.
\newblock \emph{Preprint}, 2026.

\bibitem[Spirtes et~al.(2000)]{spirtes2000causation}
Spirtes, P., Glymour, C., and Scheines, R.
\newblock \emph{Causation, Prediction, and Search}.
\newblock MIT Press, 2nd edition, 2000.

\bibitem[Wang et~al.(2023)]{wang2023interpretability}
Wang, K., Variengien, A., Conmy, A., Shlegeris, B., and Steinhardt, J.
\newblock Interpretability in the wild: A circuit for indirect object identification in {GPT}-2 small.
\newblock In \emph{ICLR}, 2023.

\bibitem[Woodward(2003)]{woodward2003making}
Woodward, J.
\newblock \emph{Making Things Happen: A Theory of Causal Explanation}.
\newblock Oxford University Press, 2003.

\end{thebibliography}

% ============================================================
\appendix

\section{IOI Claim Specification (Example)}
\label{app:ioi-spec}

The IOI claim specification defines a 6-step computational DAG:

\begin{enumerate}[nosep]
    \item \textbf{Duplicate Token Detection} (DTH: heads 0.1, 3.0) --- detects that a name appears twice.
    \item \textbf{Previous Token Tracking} (PTH: heads 2.2, 4.11) --- attends to the token before a name.
    \item \textbf{Induction} (IND: heads 5.5, 6.9) --- uses DTH and PTH signals to identify the repeated name.
    \item \textbf{S-Inhibition} (S-Inh: heads 7.3, 7.9, 8.6, 8.10) --- suppresses the repeated subject name.
    \item \textbf{Name Mover} (NM: heads 9.9, 9.6, 10.0) --- copies the non-repeated (indirect object) name to the output.
    \item \textbf{Negative Name Mover} (NegNM: heads 10.7, 11.10) --- provides opposing signal.
\end{enumerate}

\paragraph{Positive predictions (7).}
\begin{itemize}[nosep]
    \item Ablating DTH reduces output ($\geq 0.2$ decrease in logit difference).
    \item Ablating induction heads substantially reduces logit difference ($\geq 0.5$).
    \item Ablating S-Inh kills output ($\geq 0.8$ decrease).
    \item Ablating S-Inh reduces name mover activation ($\geq 0.2$).
    \item Ablating NegNM increases logit difference (removes opposing signal).
    \item Ablating DTH reduces induction head activation (tests DTH$\rightarrow$IND edge).
    \item Ablating PTH reduces induction head activation (tests PTH$\rightarrow$IND edge).
\end{itemize}

\paragraph{Negative controls (5).}
\begin{itemize}[nosep]
    \item Ablating NM does not affect upstream S-Inh.
    \item Ablating NegNM does not affect upstream S-Inh.
    \item Ablating PTH alone does not destroy circuit output ($< 0.3$ decrease).
    \item Ablating S-Inh does not affect upstream DTH.
    \item Ablating NM does not affect NegNM (parallel roles, not connected).
\end{itemize}

This specification encodes the directed causal structure of the proposed mechanism and provides concrete, falsifiable predictions that can be tested automatically by the verification pipeline.

\section{Complete Metric List by Evidence Family}
\label{app:metrics}

\paragraph{Causal (32 metrics).} activation\_patching, logit\_diff, role\_ablation, causal\_scrubbing, sigma\_ablation, resample\_complement, misalignment, das\_iia, iia\_variants, path\_patching, counterfactual\_consistency, multi\_axis\_iia, corrupt\_restore, mediation, mediation\_v2, pse, cate, intervention\_specificity, eap, atp\_star, pairwise\_synergy, shapley\_interactions, replacement\_test, composition\_test, operation\_specification, held\_out\_prediction, procedure\_specification, logic\_gates, cross\_task\_transfer, cross\_model\_invariance, minimality\_class, intermediate\_state\_prediction.

\paragraph{Structural (18 metrics).} k\_composition, copying\_score, qk\_norms, cmd, edge\_jaccard, weight\_eap\_jaccard, network\_motifs, motif\_enrichment, attention\_clustering, capacity\_utilization, k\_alignment, weight\_extended, spectral\_svd, path\_identification, edge\_necessity, path\_specificity, compositional\_sufficiency, graph\_minimality.

\paragraph{Information-theoretic (9 metrics).} pid, ocse, notears, mutual\_information, conditional\_mi, granger\_causality, info\_bottleneck, o\_information, transfer\_entropy.

\paragraph{Behavioral (20 metrics).} effect\_size, dose\_response, ce\_delta, per\_token\_nll, calibration, generalization\_gap, mdl\_compression, subnetwork\_probe, output\_variants, output\_variants\_kl, output\_variants\_topk, corrupt\_restore\_behavioral, mean\_centered\_logit, normative\_account, error\_boundary, boundary\_sweep, epistemic\_gradient, cross\_task\_generalization, hyperparam\_sensitivity, llc.

\paragraph{Representational (5 metrics).} attention\_entropy, cka, cka\_cross\_arch, probe\_decodability, causal\_representation.

\end{document}
